% =====================================================================
%  GALOIS THEORY OF THE HORIZON COVER  --  DRAFT v0.2, 2026-07-10
%  v0.2: temperature parity corrected (Theta = S*T is the odd covariant;
%        the temperature pair is Galois-conjugate, not odd); references
%        pinned (incl. author correction on arXiv:1008.0585).
%  Author: Will Lloyd (Lloyd Geometric Diagnostics)
%  Status: WORKING DRAFT. Not for distribution or submission.
%  Prior-art status: preliminary sweep complete (PRIOR_ART_REPORT_2026-07-10);
%  INSPIRE-HEP full-text sweep and expert consultation OUTSTANDING.
%  All novelty statements are provisional ("to our knowledge").
%  Title alternatives (parked):
%    - The Two Horizons Are a Galois Pair
%    - What the Area Product Knows
%    - Vieta's Formulas for Black Holes
% =====================================================================
\documentclass[11pt]{article}
\usepackage[margin=1.05in]{geometry}
\usepackage{amsmath,amssymb,amsthm}
\usepackage[colorlinks=true,linkcolor=blue,citecolor=blue,urlcolor=blue]{hyperref}
\usepackage{booktabs}

\newtheorem{theorem}{Theorem}[section]
\newtheorem{proposition}[theorem]{Proposition}
\newtheorem{lemma}[theorem]{Lemma}
\newtheorem{corollary}[theorem]{Corollary}
\theoremstyle{definition}
\newtheorem{definition}[theorem]{Definition}
\theoremstyle{remark}
\newtheorem{remark}[theorem]{Remark}
\newtheorem{observation}[theorem]{Observation}

\newcommand{\Q}{\mathbb{Q}}
\newcommand{\ZZ}{\mathbb{Z}}
\newcommand{\disc}{\Delta}
\newcommand{\Gal}{\mathrm{Gal}}
\newcommand{\Res}{\mathrm{Res}}

\title{\textbf{Galois Theory of the Horizon Cover}\\[4pt]
\large Black-hole thermodynamics as invariant theory,\\
and an Abel--Ruffini obstruction at four charges\\[6pt]
\normalsize\fbox{DRAFT v0.2 --- 10 July 2026 --- not for distribution}}
\author{Will Lloyd\\
\small Lloyd Geometric Diagnostics Pty Ltd, Coffs Harbour, NSW, Australia\\
\small \texttt{[email placeholder]}}
\date{}

\begin{document}
\maketitle

\begin{abstract}
The inner and outer horizon entropies of a Kerr--Newman black hole are the two
roots of a monic quadratic whose coefficients are \emph{rational} in the
asymptotic charges: $z^2-e_1 z+e_2=0$ with $e_1=2\pi(2M^2-Q^2)$ and
$e_2=\pi^2(4J^2+Q^4)$. We develop the thermodynamics of the horizon pair as the
Galois theory of this cover. The $\ZZ/2$ grading by the deck transformation
(the horizon swap) reproduces the known left/right thermodynamic dictionary
--- sector entropies and temperatures, sector first laws and Smarr formulas,
the universal area product, the hidden-conformal temperatures and central
charge, and extremal chirality --- from the single principle that
\emph{invariants descend to the observables and odd data does not}. The
Galois-odd datum of the pair is the covariant $\Theta=S\,T$; the temperature
itself is mixed under the horizon swap, $T_+\mapsto T_-$, and the two
temperatures form their own conjugate pair with rational symmetric functions. For the four-charge
(Cveti\v{c}--Youm) generalization the invariant persists,
$e_2=\pi^2\!\left(4J^2+N_1N_2N_3N_4\right)$ with $N_i=4Q_i$, but the cover
degree over the physical field $\Q(M,Q_1,\dots,Q_4)$ jumps to five. We certify,
by Dedekind reduction at multiple rational points, that the Galois group of the
seed-mass fiber is $S_5$; consequently no closed-form radical expression exists
for the mass fiber, and --- via an exact collapse of the radical tower --- none
exists for the horizon entropies themselves: $(S_++S_-)^2/\pi^2$ is a primitive
element of the $S_5$ field, while $S_+S_-$ remains rational. The solvability
threshold sits exactly at four charges. The Galois norm of the mass fiber
factors completely into the sixteen BPS-type linear forms
$M\pm Q_1\pm Q_2\pm Q_3\pm Q_4$, exhibiting the marginal-stability walls as
the vanishing locus of a norm. To our knowledge the algebraic-obstruction layer
of these results has not appeared in the literature; a full prior-art search is
in progress.
\end{abstract}

% =====================================================================
\section{Introduction}
% =====================================================================

Two families of exact results dominate the thermodynamics of multi-horizon
black holes. The first is the left/right (``inner mechanics'') apparatus: sector
entropies $S_{L,R}=(S_+\pm S_-)/2$, sector temperatures, first laws and Smarr
formulas on both horizons, and negative inner-horizon temperatures, developed
by Cveti\v{c} and Larsen \cite{CL1,CL2} and given a systematic Gibbsian
formulation by Cveti\v{c}, Gibbons, L\"u and Pope \cite{CGLP2018}. The second
is the universal area product: $S_+S_-$ is independent of the mass and
plausibly quantized \cite{AH,CGP2011,CR2012}, a fact used as evidence for a
conformal-field-theory description of the microstates in the hidden-conformal
program \cite{CMS,WangLiu,ChenLong4Q}.

Both families are computed, throughout the literature, by \emph{manipulating
the roots}: one writes $S_+$ and $S_-$ explicitly in an auxiliary
parametrization (the solution-generating frame: a seed mass $m$ and boosts
$\delta_i$ \cite{CY,Std4Q}) and forms the combinations of interest. This paper
starts from an elementary observation about that method:

\begin{quote}
\emph{The roots exist only in the auxiliary frame. In the physical variables
--- the asymptotic charges an observer measures --- the objects that survive
are exactly the Galois invariants of the horizon cover.}
\end{quote}

For Kerr--Newman the cover is a quadratic and the distinction is soft: both
$e_1=S_++S_-$ and $e_2=S_+S_-$ are rational in $(M,J,Q)$, and $\sqrt{\disc}$
(one square root of an invariant) recovers the roots. We show in
Section~\ref{sec:KN} that the entire classical left/right dictionary is the
representation theory of the $\ZZ/2$ deck action on this quadratic: the even
(invariant) sector is the left-moving sector, the odd sector is the
right-moving sector, and the Galois-odd datum of the pair is the covariant
$\Theta=S\,T$ --- the orientation the invariant ring discards.

At four charges the distinction becomes a theorem. The seed mass obeys
\begin{equation}
4M \;=\; \sum_{i=1}^{4}\sqrt{\,m^2+N_i^2\,},\qquad N_i=4Q_i ,
\label{eq:massfiber}
\end{equation}
a four-radical relation whose norm form is a \emph{quintic} in $u=m^2$ over
$\Q(M,Q_1,\dots,Q_4)$. We certify (Theorem~\ref{thm:S5}) that its Galois group
is $S_5$: by the Abel--Ruffini theorem there is no closed-form radical
expression for the seed mass in the physical variables. The parametric
(cosh/sinh) presentation of four-charge thermodynamics is therefore not a
convention but a necessity. We then show (Theorem~\ref{thm:entropy}) that the
obstruction propagates to the entropy itself: for the static four-charge black
hole the radical tower above the quintic collapses --- $\prod_i(w_i+m)\cdot
\prod_i(w_i-m)=\prod_i N_i^2$ is rational --- so that
$(S_++S_-)^2/\pi^2$ is a \emph{primitive element} of the $S_5$ field, while
$S_+S_-$ is rational. The horizon entropies of a four-charge black hole cannot
be written in radicals in the observables; their symmetric functions can, and
the celebrated ``universal'' formulas of the literature are precisely the
symmetric functions --- \emph{universality is Galois descent}.

The solvability threshold is sharp (Proposition~\ref{prop:ladder}): one and
two charges are rational (the two-charge seed is, explicitly, a product of the
four BPS walls), three charges give a solvable quartic, and four charges give
the first $S_5$. Finally, the Galois norm of the mass fiber factors completely
over $\Q$ into the sixteen BPS-type linear forms (Proposition~\ref{prop:walls}),
identifying the marginal-stability walls of the charge lattice with the
vanishing locus of a norm.

\paragraph{Relation to prior work.} The thermodynamic dictionary of
Section~\ref{sec:KN} is \emph{known}: the sector first laws, Smarr formulas,
negative inner temperatures, half-energy property, entropy product and
inversion appear in \cite{CGLP2018}; the hidden-conformal temperatures and
central charge for Kerr--Newman in \cite{WangLiu} (Kerr case \cite{CMS});
the universal area products in \cite{AH,CGP2011,CR2012}; the left/right
formalism originates with \cite{CL1,CL2}. Our contribution there is the
derivation principle: all of it follows from the $\ZZ/2$ Galois grading of one
rational quadratic, with no geometric or conformal input. The
number-theoretic study of \emph{extremal} black holes --- attractor arithmetic,
class numbers, complex multiplication \cite{Moore,BKO} --- concerns the
abelian Galois theory of attractor values and is complementary to the
questions here, which concern the non-abelian Galois theory of the
\emph{non-extremal cover}. To our knowledge, the algebraic-obstruction layer
--- the $S_5$ certification, the Abel--Ruffini statements for mass and entropy,
the charge-count threshold, and the wall-factorization of the norm --- has not
appeared in the literature; a full prior-art search (INSPIRE full-text and
expert consultation) is in progress, and these claims should be read with that
caveat. Notably, \cite{CGLP2018} states the entropy--mass \emph{inversion}
formula for the quadratic case and develops the sector apparatus at length
without raising the solvability question in the multi-charge case; a full-text
search of that paper for the relevant algebraic vocabulary returns no
occurrences.

\paragraph{Conventions.} Units $G=\hbar=c=1$; entropy $S=A/4$. For
Kerr--Newman we use the fundamental relation in the form
\begin{equation}
M^2 \;=\; U(S,J,Q)\;=\;\frac{S}{4\pi}+\frac{\pi J^2}{S}+\frac{Q^2}{2}
+\frac{\pi Q^4}{4S},
\label{eq:U}
\end{equation}
with conjugate potentials $T=U_S/2M$, $\Omega=U_J/2M$, $\Phi=U_Q/2M$.
All computer-verified statements are exact (symbolic or exact rational
arithmetic); the scripts, their SHA-256 hashes, and the Dedekind witnesses are
inventoried in Appendix~\ref{app:certs}.

% =====================================================================
\section{The Kerr--Newman horizon pair as a Galois cover}
\label{sec:KN}
% =====================================================================

\subsection{The quadratic, its invariants, and the horizon swap}

Solving \eqref{eq:U} for $S$ at fixed $(M,J,Q)$ gives the monic quadratic
\begin{equation}
S^2 - e_1 S + e_2 = 0,\qquad
e_1 = 2\pi\,(2M^2-Q^2),\qquad
e_2 = \pi^2\,(4J^2+Q^4),
\label{eq:quadratic}
\end{equation}
whose roots are the outer and inner horizon entropies
\begin{equation}
S_\pm \;=\; \pi\bigl(2M^2-Q^2 \pm 2\sqrt{\disc}\,\bigr),
\qquad \disc \;=\; M^4 - M^2Q^2 - J^2 ,
\end{equation}
with $e_1^2-4e_2 = 16\pi^2\disc$. Both coefficients are \emph{rational} in the
asymptotic charges; $e_2$ is independent of the mass (the universal area
product \cite{AH,CGP2011} in entropy normalization). The Galois group of
\eqref{eq:quadratic} over $\Q(M,J,Q)$ is $\ZZ/2$, realized geometrically as the
\emph{horizon swap} $S_+\leftrightarrow S_-$, i.e.\ the deck transformation
$\sqrt{\disc}\mapsto-\sqrt{\disc}$ of the double cover branched along the
extremal locus $\disc=0$.

Three elementary facts organize everything that follows (all verified
symbolically; Appendix~\ref{app:certs}):
\begin{enumerate}
\item \emph{Monodromy.} Transporting around $\disc=0$ swaps the sheets while
$e_1,e_2$ are single-valued. The branch point is simple everywhere on the
smooth extremal edge: $\disc=0$ and $\nabla\disc=0$ have no common solution
away from $M=J=Q=0$.
\item \emph{Past the wall.} For $\disc<0$ (the over-extremal region) the pair
becomes complex-conjugate, $S_-=\overline{S_+}$, with
$\mathrm{Re}\,S_\pm=e_1/2$ frozen at the extremal value and
$\mathrm{Im}\,S_\pm=\pm 2\pi\sqrt{|\disc|}$. The individual root
``goes complex''; the invariants cross the wall rationally. \emph{The entropy
never crosses; its symmetric functions do.}
\item \emph{Vieta identities for the potentials.} On the two branches,
\begin{equation}
U_S(S_+)=+\frac{\sqrt\disc}{S_+},\qquad
U_S(S_-)=-\frac{\sqrt\disc}{S_-},
\end{equation}
so that $T_\pm = \pm\sqrt{\disc}/(2MS_\pm)$: the inner-horizon temperature is
negative in the Gibbsian convention, in agreement with \cite{CGLP2018}.
\end{enumerate}

\subsection{The Galois-odd datum is $\Theta=S\,T$}

\begin{proposition}[Splitting--temperature identity]
\label{prop:splitting}
Let $\Theta:=S\,T$ on either branch, so that $\Theta_\pm=\pm\sqrt{\disc}/2M$.
Then
\[
S_+ - S_- \;=\; 4\pi\sqrt{\disc}\;=\;8\pi M\,\Theta_+ .
\]
Under the horizon swap $\Theta$ is odd, $\tau(\Theta_\pm)=-\Theta_\pm$, with
invariant square $\Theta^2=\disc/4M^2$. The temperature itself is mixed ---
$\tau(T_+)=T_-$, an odd numerator over a branch-dependent denominator --- and
the two temperatures are Galois-conjugate over the observables, with rational
symmetric functions (for $e_2\neq0$)
\[
T_+ + T_- \;=\; -\frac{2\pi\disc}{M\,e_2},\qquad
T_+ T_- \;=\; -\frac{\disc}{4M^2 e_2}.
\]
At $J=Q=0$ the identity returns the Schwarzschild value $T_+=1/8\pi M$
exactly.
\end{proposition}

The separation of the two sheets \emph{is} the odd covariant. Since
$(S_+-S_-)^2=e_1^2-4e_2$, the invariant ring sees $\Theta$ only through its
square: the sign of $\Theta$ --- outer versus inner, the arrow that
distinguishes the horizons --- is precisely the datum the Galois quotient
forgets. Temperature inherits that arrow through $T=\Theta/S$, which is how
the negative inner temperature of \cite{CGLP2018} arises, but $T$ itself is
not of pure parity: it is the odd covariant divided by a root. This is the
organizing fact of the paper: \emph{the thermodynamic quantities that are
hard to define in the observables are the odd and mixed ones; the ones that
``descend'' are the invariants.}

\subsection{The left/right dictionary from the grading}
\label{sec:dictionary}

Define the even and odd sector entropies
\begin{equation}
S_L=\tfrac12(S_++S_-)=\tfrac12 e_1=\pi(2M^2-Q^2),
\qquad
S_R=\tfrac12(S_+-S_-)=2\pi\sqrt{\disc},
\end{equation}
and sector temperatures by thermodynamic conjugacy at fixed $(J,Q)$:
$1/T_{L,R}:=\partial_M S_{L,R}$. All of the following are exact identities
(Appendix~\ref{app:certs}):
\begin{align}
T_L &= \frac{1}{4\pi M} \quad\text{(mass only: $J$ and $Q$ drop out)},&
T_R &= \frac{\sqrt{\disc}}{M\,e_1},\\
\frac{1}{T_\pm} &= \frac{1}{T_L}\pm\frac{1}{T_R}
\quad\text{(harmonic composition)},&
\Bigl(\frac{T_R}{T_L}\Bigr)^{\!2} &= 1-\frac{4e_2}{e_1^2}
= \frac{S_R^2}{S_L^2}.
\end{align}
The left/right asymmetry is the normalized discriminant of the quadratic. The
sector split, the half-energy property of the left sector, and the harmonic
composition agree with \cite{CGLP2018,CL1}; the Cardy-consistent
hidden-conformal temperatures of \cite{WangLiu},
$T_L^{\rm CFT}=(2M^2-Q^2)/4\pi J$ and
$T_R^{\rm CFT}=\sqrt{M^4-J^2-M^2Q^2}/2\pi J$ with $c_L=c_R=12J$, are
$e_1/8\pi^2 J$ and $\sqrt{\disc}/2\pi J$ in the present variables, and the
Cardy formula returns exactly $S_L$ and $S_R$ (verified symbolically; Kerr
case \cite{CMS}). Extremal chirality --- $S_R,T_R\to 0$ at $\disc\to0$, the
purely left-moving limit of \cite{GHSS} --- is, in this language, the death of
the odd sector at the branch point.

\begin{proposition}[Sector first laws and Smarr formulas; charge parity]
\label{prop:sectors}
With sector potentials $\Omega_X=-T_X\,\partial_J S_X$,
$\Phi_X=-T_X\,\partial_Q S_X$ one finds, exactly,
\begin{align}
\Omega_L &\equiv 0, &
\Phi_L &= \frac{Q}{2M}, &
M &= 2T_LS_L+\Phi_L Q,\\
\Omega_R &= \frac{J}{M(2M^2-Q^2)}, &
\Phi_R &= \frac{MQ}{2M^2-Q^2}, &
M &= 2T_RS_R+2\Omega_R J+\Phi_R Q .
\end{align}
The left sector is a complete, \emph{non-rotating}, charged thermodynamic
system: the angular momentum couples exclusively to the odd sector. Charge
behaves differently: $Q$ appears in both $e_1$ and $\disc$, and couples to
both sectors.
\end{proposition}

Sector first laws and Smarr formulas for the two-horizon system appear in
\cite{CGLP2018}; we re-derive them here from the grading and record the
following organizing observation, which to our knowledge has not been isolated
as a principle (pending the picture-literature check, cf.\ \cite{Twofold}):

\begin{observation}[Charge-parity principle]
\label{obs:parity}
A conserved charge couples to the sectors according to which symmetric
functions of the quadratic carry it. $J$ appears only in the discriminant
(equivalently only in $e_2$): it is \emph{parity-pure}, purely right-moving,
and $\Omega_L\equiv0$. $Q$ leaks into $e_1$: it \emph{straddles}. This is why
the $J$-picture of the hidden-conformal correspondence is grading-natural
(the odd sector's circle is intrinsically the rotation angle, $c=12J$ with no
further choices), while charged pictures require auxiliary structure --- the
well-known arbitrary scale in their central charges is, in this language, the
statement that no parity-pure circle exists for $Q$. On the
Reissner--Nordstr\"om slice no charge is parity-pure at all.
\end{observation}

\subsection{Attractor as Vieta; duality as the symmetry of $e_2$}

At the branch point the double root gives
$S_{\rm ext}=e_1/2=\sqrt{e_2}=\pi\sqrt{4J^2+Q^4}$: the extremal entropy is the
square root of the constant term, mass-independent by inspection --- the
attractor-mechanism entropy formula \cite{Moore} appearing as Vieta's formula
at the discriminant-zero locus. For dyonic Kerr--Newman
($Q^2\to Q_e^2+Q_m^2$) the electromagnetic duality group is exactly the
symmetry group of $e_2$: $e_1,e_2,\disc$ are $SO(2)$-invariant, the sector
potentials transform as duality doublets, both sector Smarr formulas hold
dyonically, and $T_L=1/4\pi M$ is duality-blind (all verified exactly).

% =====================================================================
\section{Four charges: the invariant persists}
\label{sec:fourcharge}
% =====================================================================

The four-charge generalization (the Cveti\v{c}--Youm class \cite{CY,Std4Q})
is parametrized by a seed mass $m$, rotation $a$, and four boosts $\delta_i$:
\begin{equation}
M=\frac{m}{4}\sum_i\cosh2\delta_i,\quad
Q_i=\frac{m}{4}\sinh2\delta_i,\quad
J=ma\,(\Pi_c-\Pi_s),\quad
N_i:=4Q_i ,
\end{equation}
\begin{equation}
S_\pm = 2\pi m\Bigl[\,m(\Pi_c+\Pi_s)\;\pm\;\sqrt{m^2-a^2}\,(\Pi_c-\Pi_s)\Bigr],
\qquad \Pi_c=\prod_i\cosh\delta_i,\ \Pi_s=\prod_i\sinh\delta_i .
\label{eq:fourchargeS}
\end{equation}
Two exact anchors pin this presentation to Section~\ref{sec:KN}: the
zero-boost limit reproduces the Kerr quadratic, and the equal-charge diagonal
reproduces the Kerr--Newman quadratic \eqref{eq:quadratic} with
$Q=\sum_iQ_i$ (both verified symbolically). The second symmetric function
generalizes exactly:

\begin{proposition}[The invariant at four charges]
\label{prop:crown}
As an unconditional algebraic identity in the parametrization,
\begin{equation}
e_2 \;=\; S_+S_- \;=\; \pi^2\bigl(4J^2 + N_1N_2N_3N_4\bigr).
\end{equation}
The Kerr--Newman $Q^4$ is the equal-charge diagonal of the four-fold product.
The mass enters only through $(J,N_i)$; with quantized charges the invariant
lies on a lattice --- the classical shadow of level matching. The static
sectors are $S_{L,R}=2\pi m^2(\Pi_c\pm\Pi_s)$: the Galois-even/odd split is the
cosh/sinh product split of the string-theoretic left/right excitation counting
\cite{CL1,SV}.
\end{proposition}

% =====================================================================
\section{The unsolvability theorems}
\label{sec:unsolvable}
% =====================================================================

\subsection{The mass fiber and its norm}

Eliminating the boosts from the parametrization yields the four-radical
relation \eqref{eq:massfiber} for the seed mass. Its norm form
\begin{equation}
\mathcal N(u)\;=\;\prod_{\varepsilon\in\{\pm1\}^4}
\Bigl(4M-\sum_i \varepsilon_i\sqrt{u+N_i^2}\Bigr),
\qquad u=m^2,
\label{eq:norm}
\end{equation}
is a polynomial in $u$ of degree exactly $5$, with leading coefficient
$-2^{24}M^6$ independent of the charges (Appendix~\ref{app:degree}): the two
aligned sign patterns contribute one $u$, the four $\pm2$-pairs contribute
four, and the six balanced $(2,2)$ patterns contribute the $M^6$.

\begin{proposition}[The charge ladder]
\label{prop:ladder}
For $n$ charges the analogous norm has $u$-degree $1,1,4,5$ at
$n=1,2,3,4$. In particular: for $n\le2$ the seed is a rational function of the
observables --- explicitly, at two charges,
\begin{equation}
m^2=\frac{4\,(M-Q_1-Q_2)(M+Q_1+Q_2)(M-Q_1+Q_2)(M+Q_1-Q_2)}{M^2},
\end{equation}
the product of the four BPS walls; at $n=3$ the seed obeys an irreducible
quartic (solvable by radicals, though we are not aware of the formula having
been written); and $n=4$ is the first quintic. The solvability threshold sits
exactly at four charges.
\end{proposition}

\subsection{Abel--Ruffini for the mass fiber}

\begin{theorem}[$S_5$ obstruction, seed mass]
\label{thm:S5}
Over $\Q(M,Q_1,\dots,Q_4)$ the polynomial $\mathcal N(u)$ is irreducible with
Galois group $S_5$. Consequently the seed mass-square $u=m^2$ of the
non-extremal four-charge black hole admits no closed-form expression in
radicals in the physical variables, and the boost parametrization cannot be
inverted in radicals.
\end{theorem}

\begin{proof}[Proof (computer-certified; Appendix~\ref{app:certs})]
By the specialization principle, the Galois group of any separable rational
specialization embeds in the generic group; hence a single specialization with
non-solvable group forces the generic group to be non-solvable, and, being a
transitive subgroup of $S_5$ witnessing $S_5$ cycle types, to equal $S_5$. At
three independent rational points --- including the integer charge vector
$M=15$, $Q=(1,2,3,5)$, chosen \emph{directly} rather than through the
parametrization, above the BPS bound $M>\sum_iQ_i$ --- the specialized quintic
is separable and irreducible over $\Q$, the physical (all-plus branch) root
lies on it, and Dedekind reduction exhibits Frobenius elements of cycle type
$(2,3)$ at $p=53$, $23$, $31$ respectively. A $(2,3)$ element has order six and
exists in no transitive quintic group except $S_5$. All five roots are real at
each point (the quintics are totally real): the physical seed has four real
Galois-conjugate siblings over the same observables.
\end{proof}

\subsection{Abel--Ruffini for the entropy}

The entropy involves the boosts as well as the seed; a priori the obstruction
of Theorem~\ref{thm:S5} might be bypassed. It is not: the tower collapses.

\begin{lemma}[Collapse of the radical tower; static case]
\label{lem:collapse}
Let $w_i=\sqrt{m^2+N_i^2}$ and $u=m^2$. Then, exactly:
\emph{(i)} the branch-compatible signed radicals lie in the quintic field ---
there exist $g_i\in\Q[u]/(\mathcal N)$ with $g_i^2\equiv u+N_i^2$ and
$\sum_i g_i\equiv 4M$ (certified as polynomial identities modulo
$\mathcal N$); \emph{(ii)} with $A=\prod_i(w_i+m)$, $B=\prod_i(w_i-m)$ one has
$AB=\prod_iN_i^2\in\Q$, so that $y:=(S_++S_-)/\pi=\sqrt A+\sqrt B$ satisfies
\begin{equation}
y^2 \;=\; \gamma(u)\;:=\;2\Bigl[e_4(w)+u\,e_2(w)+u^2+\textstyle\prod_iN_i\Bigr]
\;\in\;\Q[u]/(\mathcal N).
\end{equation}
\end{lemma}

\begin{theorem}[Entropy non-radical; static four-charge]
\label{thm:entropy}
$\gamma$ is a primitive element of the quintic field: its resolvent
$\Res_u(\mathcal N,\,X-\gamma)$ is an irreducible quintic whose Galois group
is again $S_5$ (Dedekind witnesses of type $(2,3)$ at $p=17,19,23$). Since
$y^2=\gamma$, the field $\Q(M,\vec Q)(S_++S_-)$ contains the quintic field and
its Galois closure surjects onto $S_5$: it is non-solvable. Hence, in the
physical variables:
\begin{enumerate}
\item $S_++S_-$ admits no radical expression; its exact degree is $10$
(the polynomial $\Res_u(\mathcal N,\,Y^2-\gamma)$ is irreducible);
\item $S_+S_-=\pi^2\prod_iN_i$ is rational, so $S_+$ and $S_-$ individually
admit no radical expression, and neither do $S_L$ and $S_R$;
\item $S_L^2=(\pi^2/4)\,\gamma$: the left-sector level count is itself a
primitive element of the unsolvable field, while the level-matched
combination $S_L^2-S_R^2=e_2$ is rational.
\end{enumerate}
\end{theorem}

\begin{remark}[Scope]
The rotating case adds $J$ to the same tower over the same quintic and is a
bookkeeping extension, deferred. The non-radicality of the sector temperatures
$T_{L,R}$ is plausible by inheritance but not certified here. The five sheets
of the quintic are realized, at the certified points, by the all-plus sign
pattern and the four single-flip patterns: the Galois orbit of a four-charge
black hole over its own observables consists of the hole and its four
one-charge-conjugate shadows. Single charge conjugations act on the base
$(e_1,e_2)$ data --- they move between quintic sheets --- while the deck
transformation of each entropy pair remains
$\sqrt{\disc}\mapsto-\sqrt{\disc}$.
\end{remark}

\begin{remark}[Methodological corollary]
The left/right apparatus of \cite{CL1,CL2,CGLP2018} is constructed by explicit
manipulation of $S_\pm$ in the boost frame. Theorems~\ref{thm:S5} and
\ref{thm:entropy} show that at four charges this construction cannot be
transported to the observables even in principle: the roots, the sector
entropies, and the individual level counts are not radical functions of
$(M,\vec Q)$. The quantities of the apparatus that are expressible in the
observables are precisely the Galois invariants --- the area product, the
level-matched difference, the wall data of Section~\ref{sec:walls}.
\emph{Universality is descent: universal formulas are universal because they
are norms.}
\end{remark}

% =====================================================================
\section{The norm and the walls}
\label{sec:walls}
% =====================================================================

\begin{proposition}[Wall factorization of the norm]
\label{prop:walls}
The product of the five branches of the mass fiber is, exactly,
\begin{equation}
\prod_{k=1}^{5}u_k \;=\;
\frac{256}{M^{6}}\;
\prod_{\varepsilon\in\{\pm1\}^4}\bigl(M+\varepsilon_1Q_1+\varepsilon_2Q_2
+\varepsilon_3Q_3+\varepsilon_4Q_4\bigr),
\end{equation}
with leading coefficient $\mathrm{lc}(\mathcal N)=-2^{24}M^6$ independent of
the charges. Equivalently $\mathcal N(0)=4^{16}\prod_{16}(M\pm Q_1\pm Q_2\pm
Q_3\pm Q_4)$: the constant term of the quintic is the black hole evaluated at
its own BPS limit ($m\to0$), and the Galois norm vanishes exactly on the union
of the sixteen marginal-stability walls. The norm is positive throughout the
physical chamber $M>\sum_i|Q_i|$ (consistently with the totally real quintics
of Theorem~\ref{thm:S5}) and changes sign across walls: the chamber structure
of the charge lattice is read off one coefficient.
\end{proposition}

Thus the invariant tower is, floor by floor, \emph{more} rational than
solvability requires: the roots need $S_5$; the quadratic invariant
$e_2=\pi^2(4J^2+\prod N_i)$ is rational and quantized; and the quintic's top
invariant does not merely remain rational --- it splits completely over $\Q$
into the BPS walls. What the Galois quotient keeps is not just rational; it is
rational \emph{and physical}. At two charges the wall product \emph{is} the
seed (Proposition~\ref{prop:ladder}); at four charges it survives only as the
norm of five entangled sheets.

% =====================================================================
\section{Discussion and open problems}
\label{sec:discussion}
% =====================================================================

The picture that emerges is a division of black-hole thermodynamics into an
invariant layer and an odd/root layer, separated --- at four charges --- by a
non-solvable wall. The invariant layer (area products, level matching, wall
data, the left sector's half-energy) descends to the observables and is where
the literature's ``universal'' results live; the root layer (individual
entropies, individual level counts, the sign of the temperature) exists only
in the duality frame, and the solution-generating map is, in a precise sense,
radical-one-way. The classical dictionary of the hidden-conformal program is
reproduced by the grading alone, with no conformal input; the classical
checks therefore carry no discriminating power about the existence of a dual
CFT, and the burden of evidence rests on the genuinely quantum observables.

Open problems, in rough order of proximity: \emph{(i)} the Galois/monodromy
group of the full four-charge entropy cover (the degree-$10$ extension and its
interplay with the $(\ZZ/2)^4$ charge conjugations), for which imprimitive
wreath-type methods are the natural tool; \emph{(ii)} the rotating extension
of Theorem~\ref{thm:entropy} and certification for $T_{L,R}$; \emph{(iii)}
wall structure, if any, of the lower symmetric functions $e_1,\dots,e_4$ of
the quintic; \emph{(iv)} the lattice interpretation of $e_2$ in string units;
\emph{(v)} the charged-picture central charges in the light of
Observation~\ref{obs:parity}; \emph{(vi)} five-dimensional and
cosmological-constant generalizations, where the horizon polynomials have
different degrees and the ladder of Proposition~\ref{prop:ladder} predicts a
different solvability frontier.

\paragraph{Acknowledgments.}
[Placeholder.]
% The author thanks a very patient artificial colleague for verification,
% adversarial review, and for refusing to let Hawking radiation be claimed
% before breakfast.

% =====================================================================
\appendix
\section{Computational certificates}
\label{app:certs}
% =====================================================================

All results marked as verified are exact: symbolic identities in
\texttt{sympy}, or exact rational arithmetic with numerics used only as
scaffolding (every numerically-guided step is discharged by an exact
polynomial identity, e.g.\ modulo $\mathcal N$). Scripts (SHA-256 prefixes):

\begin{center}
\begin{tabular}{@{}lll@{}}
\toprule
Script & Contents & SHA-256 (prefix)\\
\midrule
\texttt{four\_charge\_crown.py} &
Prop.~\ref{prop:crown}; anchors; extremal chirality &
\texttt{8ff83aef7bc0}\\
\texttt{abel\_ruffini\_black\_holes.py} &
Thm.~\ref{thm:S5}; ladder points; witnesses &
\texttt{675b3ced84f0}\\
\texttt{entropy\_sum\_resolvent.py} &
Lem.~\ref{lem:collapse}; Thm.~\ref{thm:entropy}; degree 10 &
\texttt{7f778cb81dce}\\
\bottomrule
\end{tabular}
\end{center}

\noindent Kerr--Newman identities of Section~\ref{sec:KN} (Vieta identities,
splitting--temperature, sector laws, harmonic and ratio laws, Cardy match,
attractor, duality) were verified in-session as standalone symbolic runs; the
wall factorization of Proposition~\ref{prop:walls} was verified exactly at
four rational points together with the charge-independence of the leading
coefficient.

\paragraph{Dedekind witnesses (Theorem~\ref{thm:S5}).}
Point A: $M=3$, $Q=(\frac13,\frac12,\frac25,\frac37)$; irreducible; witness
$p=53$, type $(2,3)$; $5$-cycles at $p=41,43$. Point B: $M=15$, $Q=(1,2,3,5)$;
irreducible; witness $p=23$; $5$-cycles at $p=29,43,47$. Point C: $M=2$,
$Q=(\frac27,\frac59,\frac14,\frac4{11})$; irreducible; witness $p=31$.
All quintics separable and totally real; the physical root verified on each to
residuals $\le10^{-15}$ (exact irreducibility making the identification
rigorous).

\paragraph{Witnesses (Theorem~\ref{thm:entropy}, point B).}
Exact certificates modulo $\mathcal N$: $g_i^2\equiv u+N_i^2$,
$\sum g_i\equiv 4M$, $\alpha^2-u\beta^2\equiv(\prod N_i)^2$. Resolvent
irreducible; witnesses $p=17,19,23$, all type $(2,3)$; $E(Y)$ of degree $10$
irreducible over $\Q$; physical cross-check exact to $80$ digits. Branch
patterns of the five sheets: $(++++)$ and the four single-flips.

\section{Degree count and the two-charge closed form}
\label{app:degree}

For large $u$ each factor of \eqref{eq:norm} behaves as
$4M-(\sum_i\varepsilon_i)\sqrt u+O(u^{-1/2})$. Pairing $\varepsilon$ with
$-\varepsilon$: the aligned pair contributes $16M^2-16u+\dots$, each of the
four $|\sum\varepsilon|=2$ pairs contributes $16M^2-4u+\dots$, and each of the
three balanced pairs contributes $16M^2+O(1)$; half-integer powers cancel
because $\mathcal N$ is a polynomial. Hence $\deg_u\mathcal N=1+4+0=5$ and
$\mathrm{lc}=(-16)(-4)^4(4M)^6=-2^{24}M^6$. The same count gives degrees
$1,1,4$ at $n=1,2,3$ charges; at $n=2$ the norm is linear in $u$ and solving it
yields the wall-product closed form of Proposition~\ref{prop:ladder}.

% =====================================================================
\begin{thebibliography}{99}
% =====================================================================
\bibitem{CL1} M.~Cveti\v{c} and F.~Larsen,
``General rotating black holes in string theory: Greybody factors and event
horizons,'' Phys.\ Rev.\ D {\bf 56} (1997) 4994, [hep-th/9705192].
\bibitem{CL2} M.~Cveti\v{c} and F.~Larsen,
``Greybody factors for rotating black holes in four dimensions,''
Nucl.\ Phys.\ B {\bf 506} (1997) 107, [hep-th/9706071].
\bibitem{CGLP2018} M.~Cveti\v{c}, G.\,W.~Gibbons, H.~L\"u and C.\,N.~Pope,
``Killing Horizons: Negative Temperatures and Entropy Super-Additivity,''
Phys.\ Rev.\ D {\bf 98} (2018) 106015, [arXiv:1806.11134].
\bibitem{AH} M.~Ansorg and J.~Hennig,
``The inner Cauchy horizon of axisymmetric and stationary black holes with
surrounding matter,'' Class.\ Quant.\ Grav.\ {\bf 25} (2008) 222001,
[arXiv:0810.3998]; ``Inner Cauchy horizon of axisymmetric and stationary
black holes with surrounding matter in Einstein--Maxwell theory,''
Phys.\ Rev.\ Lett.\ {\bf 102} (2009) 221102, [arXiv:0903.5405].
\bibitem{CGP2011} M.~Cveti\v{c}, G.\,W.~Gibbons and C.\,N.~Pope,
``Universal Area Product Formulae for Rotating and Charged Black Holes in
Four and Higher Dimensions,'' Phys.\ Rev.\ Lett.\ {\bf 106} (2011) 121301,
[arXiv:1011.0008].
\bibitem{CR2012} A.~Castro and M.\,J.~Rodriguez,
``Universal properties and the first law of black hole inner mechanics,''
Phys.\ Rev.\ D {\bf 86} (2012) 024008, [arXiv:1204.1284].
\bibitem{CMS} A.~Castro, A.~Maloney and A.~Strominger,
``Hidden Conformal Symmetry of the Kerr Black Hole,''
Phys.\ Rev.\ D {\bf 82} (2010) 024008, [arXiv:1004.0996].
\bibitem{WangLiu} Y.-Q.~Wang and Y.-X.~Liu,
``Hidden Conformal Symmetry of the Kerr-Newman Black Hole,''
JHEP {\bf 08} (2010) 087, [arXiv:1004.4661].
\bibitem{ChenLong4Q} K.-N.~Shao and Z.~Zhang,
``Hidden Conformal Symmetry of Rotating Black Hole with four Charges,''
Phys.\ Rev.\ D {\bf 83} (2011) 106008, [arXiv:1008.0585].
\bibitem{GHSS} M.~Guica, T.~Hartman, W.~Song and A.~Strominger,
``The Kerr/CFT Correspondence,''
Phys.\ Rev.\ D {\bf 80} (2009) 124008, [arXiv:0809.4266].
\bibitem{CY} M.~Cveti\v{c} and D.~Youm,
``Entropy of non-extreme charged rotating black holes in string theory,''
Phys.\ Rev.\ D {\bf 54} (1996) 2612, [hep-th/9603147].
\bibitem{Std4Q} Z.-W.~Chong, M.~Cveti\v{c}, H.~L\"u and C.\,N.~Pope,
``Charged rotating black holes in four-dimensional gauged and ungauged
supergravities,'' Nucl.\ Phys.\ B {\bf 717} (2005) 246, [hep-th/0411045];
the parametrization in the form used here is also displayed in
M.~Cveti\v{c}, G.\,W.~Gibbons, C.\,N.~Pope and B.\,F.~Whiting,
``Positive energy functional for massless scalars in rotating black hole
backgrounds of maximal ungauged supergravity,'' [arXiv:1912.08988].
\bibitem{SV} A.~Strominger and C.~Vafa,
``Microscopic origin of the Bekenstein-Hawking entropy,''
Phys.\ Lett.\ B {\bf 379} (1996) 99, [hep-th/9601029].
\bibitem{Moore} G.\,W.~Moore, ``Arithmetic and Attractors,''
[hep-th/9807087].
\bibitem{BKO} N.~Benjamin, S.~Kachru, K.~Ono and L.~Rolen,
``Black holes and class groups,'' Res.\ Math.\ Sci.\ (2018),
[arXiv:1807.00797]; M.~G\"unaydin, S.~Kachru and A.~Tripathy,
``Black holes and Bhargava's invariant theory,'' [arXiv:1903.02323].
\bibitem{Twofold} C.-M.~Chen, Y.-M.~Huang, J.-R.~Sun, M.-F.~Wu and S.-J.~Zou,
``Twofold Hidden Conformal Symmetries of the Kerr-Newman Black Hole,''
Phys.\ Rev.\ D {\bf 82} (2010) 066004, [arXiv:1006.4097].
  %% NOTE: unswept for the charge-parity principle; check before claiming
  %% Observation 2.8 as new.
\bibitem{Dummit} D.\,S.~Dummit, ``Solving solvable quintics,''
Math.\ Comp.\ {\bf 57} (1991) 387.
\end{thebibliography}

\end{document}
